% =====================================================================
% draft_sections/framework_results.tex
% Results subsection: MRMA on the primary backbone + honest transfer study.
% Every number recomputed from audit.json via gen_ablation_honest.py
% (FULL-vs-ONES ablation) and verified against the audit; nothing hardcoded.
%   \ref{tab:framework_main}     -- instantiations vs baselines (raw MSE per cell)
%   \ref{tab:framework_backbone} -- FULL vs ONES mask ablation per backbone
% =====================================================================

\subsection{The Mechanism, and How Far It Transfers}
\label{sec:frameworkresults}

The central empirical question is not whether one tuned model wins, but what the
MRMA masks contribute \emph{in isolation}, and whether that contribution
transfers when the same masks are dropped into a different channel-mixing
backbone. Table~\ref{tab:framework_backbone} answers both with a single controlled
comparison: for each backbone we hold the model, its parameters and the training
recipe fixed, and replace the three relation masks with an all-ones
(fully-connected) mask (\textsc{Ones}). The mask buffer is then the \emph{only}
difference, so any change in error is causally attributable to the injected
structure.

\textbf{On the primary backbone the masks help everywhere.}
\mr{}${=}$\itrans{}${+}$MRMA beats its own all-ones control in \emph{all 16 of 16}
PEMS cells, by a mean of $+9.7\%$ and up to $+17.9\%$, with the margin growing at
longer horizons where spatial propagation dominates short-range autocorrelation.
Against the \emph{unmodified} \itrans{} backbone the same instantiation wins all
16 cells as well, up to $+21.5\%$ at PEMS04 pred-96. Because the all-ones control
keeps every parameter and only removes the structural content, this pins the gain
on the specific relation graphs rather than on sparsity or extra capacity:
sparsity applied without the precomputed structure is the weakest variant in
every cell.

\textbf{The benefit is backbone-dependent---a boundary we report, not hide.}
We carry the identical masks into two further backbones. On CrossFormer's
two-stage cross-dimension attention (\textsc{SECrossformer}) the masks lower error
in only 5 of 16 PEMS cells (mean $-4.9\%$ relative to its own all-ones control):
CrossFormer already factorises channel interaction into a routed two-stage
attention, so a static pre-softmax mask has little left to constrain and can even
conflict with the learned routing. On a GRU-based state-space forecaster
(\textsc{StructMamba}), whose gated recurrence mixes channels selectively rather
than through a single global attention matrix, the mechanism is neutral. The
pattern is consistent and mechanistic: precomputed structural masks pay off
exactly where a \emph{single} global channel-attention would otherwise spread
gradient uniformly across the $N(N{-}1)$ directed pairs---most of which our own
analysis (§\ref{sec:prelim}) shows to be uninformative---and add little where the
backbone already imposes its own channel structure.

\textbf{Standing against ten baselines.}
Table~\ref{tab:framework_main} places the instantiations alongside ten baselines
(DLinear~\cite{zeng2023transformers}, PatchTST~\cite{nie2022time},
Autoformer~\cite{wu2021autoformer}, FEDformer~\cite{zhou2022fedformer},
Stationary~\cite{liu2022non}, TimesNet~\cite{wu2023timesnet},
\textsc{CrossFormer}~\cite{zhang2023crossformer}, \itrans{}~\cite{liu2023itransformer},
and two further state-space baselines), all re-run under the shared protocol.
\mr{} is competitive with the strongest baseline across the grid and best or
tied-best on the largest network (PEMS07, $N{=}883$), where full cross-channel
attention is most prone to the gradient dilution we analyse. We do not claim
uniform state-of-the-art in every cell; the contribution is a zero-parameter
mechanism with a characterised regime of effectiveness, not a leaderboard sweep.
